\section{Evaluation Protocol}
\label{sec:eval}

% Keystone for the D\&B track. All counts anchored to the 2026-04-28 manifest
% (3,986,738 ok rows, 100\% read success).

\subsection{Tasks}
\label{sec:eval-tasks}

\begin{itemize}
  \item \textbf{T1 -- Per-atom partial charge regression.} Multi-target:
    \revised{ADCH, Becke, CM5, Hirshfeld, L\"owdin and Mulliken (six schemes)}
    jointly per atom. Headline:
    per-method, per-element MAE and Pearson $r$ on the test split.
  \item \textbf{T2 -- BCP property regression.} Targets: $\rho$,
    $\nabla^2\rho$, ellipticity at every QTAIM bond critical point.
    Headline: per-element-pair MAE on $\rho$, Pearson $r$ per pair,
    stratified by bond type (single, double, aromatic, metal-ligand)
    inferred from QTAIM topology.
  \item \textbf{T3 -- Bond classification.} Binary: does a QTAIM BCP
    exist between an atom pair within $2\times$ sum-of-covalent-radii.
    Compared against Mayer(Orca) / L\"owdin(ORCA) / fuzzy agreement. Headline: Macro-F1 plus
    per-element-pair confusion stats.
\end{itemize}

\subsection{Main Split}
\label{sec:eval-main-split}

The corpus is partitioned by molecular composition: \revised{each structure's
composition string -- the electronegativity-ordered formula produced by
\texttt{pymatgen}'s \texttt{Composition.formula} -- is hashed with SHA-256, mapped
to $[0,1)$, and assigned to train / val / test by cumulative thresholds
$0.8 / 0.1 / 0.1$. Because the key is an exact function of composition, every
structure sharing a composition receives the same assignment.}
Held-out sets (\S\ref{sec:eval-holdouts}) are excluded prior to \revised{hashing}, so no
structure appears as both an in-distribution test point and a stress-test point.
Hashing on composition rather than on row identity prevents near-duplicate
conformers of the same formula from crossing the train/val/test boundary.
\revised{Realized sizes are 3{,}125{,}019 / 398{,}982 / 403{,}951
(train / val / test) after excluding 58{,}786 held-out structures, which together
account for the full 3{,}986{,}738-structure corpus. Realized ratios are
79.6\,\% / 10.2\,\% / 10.3\,\% of the post-holdout corpus, within 0.4 percentage
points of the nominal targets.}
Verticals are not sampled uniformly because the split is global by formula; a
per-vertical post-split distribution table is provided in
\revised{App.~\ref{app:post_split}} as a transparency artifact.

\begingroup\color{revcolor}
\paragraph{Split leakage audit.}
Composition hashing removes exact-composition leakage but does not by itself
control scaffold- or homologue-level similarity, so we quantify both. No test
structure shares its exact composition with any training structure (0.0\,\%),
against 80.9\,\% for a row-level random split of the same corpus: same-composition
and conformer leakage are eliminated rather than reduced. Residual similarity is
real but bounded. Taking each test composition's nearest training neighbour among
compositions with an identical element set, 96.0\,\% of test compositions have such
a neighbour and it differs by a median of 5.3\,\% of the atom count. For
calibration, withholding the \texttt{tm\_react} vertical entirely raises that
median to 16.5\,\%, so the composition split sits roughly one third of the way from
a random split to a full vertical holdout. A direct homologue probe finds that
17.3\,\% of test compositions have a training composition differing by exactly one
CH$_2$ unit. We report these measurements rather than adding a Murcko scaffold
split because scaffold decomposition requires a sanitized covalent bond graph,
which is undefined for the transition-metal, ionic, multi-fragment and
reaction-path verticals that motivate this corpus: 27.0\,\% of structures contain
an element outside the organic subset that scaffold perception assumes, and
19.0\,\% contain a transition metal, lanthanide or actinide. A further observation
supports the stress-suite design in \S\ref{sec:eval-holdouts} over per-vertical
holdout: withholding \texttt{droplet} or \texttt{rna} entirely still leaves
67--68\,\% of their compositions matched exactly in training, because solvent
clusters and nucleotide fragments recur across verticals.
\endgroup

\subsection{Held-Out Evaluation Sets}
\label{sec:eval-holdouts}

Held-out sets are constructed independently of the main split. Membership is
decided from the manifest (and from \texttt{bond.lmdb} for the two metal-ligand
sets) and removed from train/val/test so no structure is evaluated as both
in-distribution test and stress-test. \revised{Where a set is subsampled
(H1, H3, H6), sampling is composition-stratified using a
BLAKE2b~\cite{10.1007/978-3-642-38980-1_8} rank over
\texttt{formula\_hill}, so the subsample is deterministic and reproducible.} Overlaps between held-out sets are allowed
(a structure can be both H7-large and H8-high-charge), but each set is generated
separately so error attributes to one axis. The headline signal per set is the
gap $\Delta = \mathrm{metric}_{\mathrm{held\text{-}out}} -
\mathrm{metric}_{\mathrm{main\,test}}$, reported alongside the absolute number;
a small $\Delta$ on H7 suggests size extrapolates, a large $\Delta$ on H1
indicates pair-level memorization. Three of the five sets (H1, H6, H7) inherit
the hardest axes from the OMol25 paper directly.

\paragraph{Sets and sizes.}
The five held-out sets together \revised{contain 60{,}578 records spanning 58{,}786
distinct structures} ($\sim$1.5\% of the dataset; \revised{overlaps allowed, and
1{,}792 records are structures belonging to two sets}):

\begin{itemize}
  \item \textbf{H1 -- Metal-ligand pairs (\revised{15{,}025}):} eleven (TM, partner) bond pairs sampled stratified across log-spaced frequency bands; selection driven by \emph{true} bonds (Mayer/fuzzy $\geq 0.3$) from \texttt{bond.lmdb}, not element co-occurrence (App.~\ref{app:heldout-h1}).
  \item \textbf{H3 -- Reactivity (\revised{12{,}506}):} composition-stratified subsample of \texttt{tm\_react} (5k), \texttt{electrolytes\_reactivity} (5k), \texttt{pmechdb} (2.5k); tests non-equilibrium and transition-state-adjacent geometries.
  \item \textbf{H6 -- Lanthanide-ligand pairs (2{,}589):} same construction as H1 applied to (Ln, partner) bonds (App.~\ref{app:heldout-h6}).
  \item \textbf{H7 -- Large systems (\revised{18{,}096}):} \texttt{n\_atoms > 250}; size extrapolation where fuzzy integration and grid-based QTAIM are most expensive and per-atom DFT SNR is hardest to maintain.
  \item \textbf{H8 -- Large net charges (\revised{12{,}362}):} \texttt{net\_charge\_abs > 4}; the high-charge tail where partial-charge schemes disagree most and SCF is most sensitive to functional and basis.
\end{itemize}

\revised{Counts are the records present in the released held-out LMDBs, verified
against each database's stored length. They fall a few records short of the
selection targets (H1 15{,}030, H3 12{,}507, H7 18{,}200, H8 12{,}393) where a
selected structure lacked a complete descriptor set at pull time.}

We deliberately use pair-level held-outs (H1, H6) rather than
whole-element held-outs because the goal of this benchmark is to produce
models that are usable across the periodic table. Holding out specific (metal, ligand) pairs keeps every element
represented in training while still measuring rare-chemistry transfer,
identical to OMol25.


\subsection{Methodology and Metrics}
\label{sec:eval-method}

All hyperparameter tuning and early stopping use the validation bucket; the test
bucket is reserved for a single final evaluation, with the same one-shot rule on
each held-out set. Metrics are MAE, RMSE, and Pearson $r$ per task. Stratifications:
\textbf{T1} per-method, per-element, per-vertical, per-charge-bucket; \textbf{T2}
per-element-pair, per-bond-type (covalent / metal-ligand / H-bond), per-vertical;
\textbf{T3} macro-F1 plus per-element-pair confusion.
